\chapter{Introduction}
\label{ch:introduction}

\section{Background and Motivation}
\label{sec:intro-background}

Federated learning (FL) lets a set of clients -- edge devices, organisations, or other data holders -- collaboratively train a shared machine learning model without any client exposing its raw data to the others or to a central server \citep{fedavg}. Each round, a central server broadcasts the current global model, every participating client trains it further on its own local data, and the server combines the returned updates -- typically by Federated Averaging (FedAvg) -- into a new global model. This is what makes FL attractive for privacy-sensitive applications: the server only ever sees model updates, never the data that produced them.

That same property is also FL's central weakness. Because the server cannot inspect a client's local data, it has no independent way to check that a returned update reflects honest training. A client under an attacker's control can submit a manipulated update instead, and if enough poisoned updates are averaged into the global model, the model itself is compromised -- either its overall accuracy degrades, or a narrower vulnerability is introduced while overall accuracy stays high enough to go unnoticed \citep{bagdasaryan2020,bhagoji2019}. This is known as \emph{model poisoning}, and defending against it while preserving FL's core privacy property -- never inspecting a client's raw data or gradients in detail -- is an open, actively studied problem.

The dominant defense strategy in the literature treats every client's update as a point in a high-dimensional vector space and looks for statistical outliers. Coordinate-wise median and trimmed-mean aggregation bound the influence any single coordinate can have \citep{byzantine}. Krum and Multi-Krum keep only the updates that sit closest to their neighbours under a Euclidean-distance criterion \citep{krum}. DnC projects updates onto their dominant spectral direction and removes the ones that stick out furthest along it \citep{dnc}. FLTrust bootstraps a small, trusted reference update from a clean root dataset held by the server and rescales client updates toward it by cosine similarity \citep{fltrust}. FoolsGold instead looks for clients whose update histories resemble each other -- the signature of colluding Sybils \citep{foolsgold}. DeFL departs from these single-round, whole-update statistics and instead tracks how a model's per-layer gradient-norm profile evolves across rounds \citep{defl}. What every one of these approaches has in common is that the notion of "what a poisoned update looks like" is fixed once, at design time, by a human researcher.

An attacker does not have to respect that fixed notion. Prior work has shown that an attacker aware of a specific deployed defense can construct updates that satisfy its statistical assumptions while still corrupting the aggregate \citep{fang2020,bhagoji2019}. In principle, this sets up an arms race between attack and defense. In practice, that arms race has historically been fought by hand: a researcher designs an attack against a defense, another researcher designs a defense against that attack, and the cycle repeats on the timescale of publications, not training rounds. Between human interventions, both sides are frozen.

Large language models (LLMs), meanwhile, have become capable enough -- and cheap enough to fine-tune with parameter-efficient methods such as LoRA \citep{lora2022} and reinforcement learning algorithms such as GRPO \citep{grpo} -- that it is now practical to ask whether an attack strategy can be \emph{learned} directly by a policy interacting with an FL environment, rather than designed by a person. This report investigates that question specifically for the attacker side of the FL security arms race.

\section{Problem Statement}
\label{sec:intro-problem}

Existing FL poisoning attacks in the literature are, almost without exception, hand-designed: a researcher specifies the perturbation (scale a layer, flip signs, add noise) and the targeting logic (which clients, how much) as a fixed procedure, sometimes tuned against one specific defense. This has two consequences relevant to this project. First, evaluating a new defense against such attacks only tells us how that defense performs against the attacks a researcher happened to think of, not against the space of attacks an adaptive adversary could discover. Second, comparing defenses against each other typically means comparing them against different, non-uniform attacks reported in different papers, rather than the identical attack under identical conditions.

This report addresses both consequences by building an attacker whose strategy -- which clients to poison, how many, and what to do to their weights -- is not specified in advance but is instead learned through reinforcement learning against a verifiable reward, and by evaluating that trained attacker against several defenses under an identical, replayed attack so that the defense is the only variable in the comparison.

\section{Research Aim and Objectives}
\label{sec:intro-objectives}

\textbf{Aim.} To design, implement, and evaluate a zero-touch, LLM-directed attacker for federated learning -- one that learns its poisoning strategy through reinforcement learning rather than a hand-authored attack template -- and to determine whether the resulting strategy generalises to classical, non-learned Byzantine-robust defenses it never trained against.

\textbf{Objectives.}
\begin{enumerate}
    \item Design an attacker agent contract in which an LLM policy selects, each round, which of a fixed pool of controllable clients to poison and composes a plan from a small set of primitive weight operators, using only statistics of its own controllable clients' honest updates.
    \item Formulate an exactly verifiable, continuous reward for the attacker, computed from the ground-truth poisoned set and the resulting model accuracy, that credits accuracy degradation and evasion while discouraging wasteful or easily-detected multi-client use.
    \item Train the attacker online with Group-Relative Policy Optimization (GRPO) against a co-trained defender LLM, using a training schedule -- freeze-and-alternate scheduling, stochastic opponent scoring, and an opponent league -- designed to keep both agents improving rather than stalling or cycling.
    \item Build a fixed-attack, varied-defense benchmark harness that replays the trained attacker's committed actions, unmodified, against a panel of independent defenses (FedAvg, an Oracle upper bound, FLTrust, Multi-Krum, DnC, DeFL, and the co-trained LLM defender).
    \item Empirically evaluate the trained attacker against this panel and analyse the results, in particular whether a defense's per-round detection accuracy (TPR/FPR/F1) predicts its realised protection of the model under attack.
\end{enumerate}

\section{Scope and Assumptions}
\label{sec:intro-scope}

This report evaluates the \emph{untargeted} attack goal -- degrading the global model's overall test accuracy -- on the MNIST dataset with a compact, 970-parameter multilayer perceptron (MLP) as the global model, under a non-IID client data partition. The attacker is modelled as a partial insider controlling a fixed minority of clients (5 of 20), consistent with a realistic compromise scenario in which the majority of participants remain honest. The reported empirical results are from a single 50-round benchmark run at one poison budget and one target accuracy-drop setting; they are explicitly presented as preliminary, pending replication across seeds, budgets, and settings (Chapter~\ref{ch:results}). Targeted (class-specific/backdoor) attacks, Sybil client creation, and collusion with a compromised aggregation server are out of scope for this report and are discussed as future work.

\section{Contribution}
\label{sec:intro-contribution}

This report makes the following contributions:
\begin{itemize}
    \item A zero-touch FL attacker whose client selection and weight-operator composition are generated directly by an LLM policy, using a general-purpose primitive-operator DSL, without prescribing any attack strategy in advance.
    \item A verifiable, continuous reward that jointly credits accuracy degradation and evasion while penalising malformed output, unnecessary client use, and non-diverse multi-client collusion.
    \item A GRPO training protocol -- freeze-and-alternate scheduling, stochastic opponent scoring, and an opponent league -- that keeps the attacker's learning signal alive against a co-evolving defender and formalises the interaction as a Stackelberg game.
    \item A fixed-attack, varied-defense benchmark harness that isolates the defense as the only variable when comparing FedAvg, an Oracle upper bound, FLTrust, Multi-Krum, DnC, DeFL, and a co-trained LLM defender against the same trained attacker.
    \item An empirical result showing that a defense's per-round detection metrics do not reliably predict its realised protection of the model, evidenced most sharply by DeFL recording the highest detection rate of any non-oracle defense in the panel while also suffering the largest accuracy collapse among the classical defenses.
\end{itemize}

\section{Report Structure}
\label{sec:intro-structure}

The remainder of this report is organised as follows. Chapter~\ref{ch:literature} reviews the relevant literature: the FL poisoning attack surface, Byzantine-robust and statistical defenses, adaptive and learning-based attacks, and the reinforcement-learning methods used to fine-tune LLM policies. Chapter~\ref{ch:methodology} describes the methodology: the FL simulation environment, the threat model, the attacker and defender agent designs, the GRPO training methodology and reward formulation, and the benchmark evaluation methodology. Chapter~\ref{ch:results} reports the results of a preliminary benchmark run, discusses what they reveal about the relationship between detection accuracy and realised protection, states the limitations of the current results, and concludes with a summary of findings and directions for future work. Appendix~\ref{app:prompts} reproduces the attacker and defender prompt contracts and the full operator DSL reference, and Appendix~\ref{app:config} reproduces the configuration and hyperparameters used to produce the reported results.
